\section{Related Work}

\subsection{Edge AI Frameworks}

Several frameworks target edge AI deployment, each with distinct design philosophies and trade-offs.

\textbf{TensorFlow Lite} \cite{tflite} provides a lightweight runtime for mobile and embedded devices, supporting post-training quantization and a converter from TensorFlow models. However, it lacks integrated training tools and hardware abstraction beyond ARM CPU and GPU backends.

\textbf{TensorFlow Lite Micro} \cite{tflitemicro} extends TensorFlow Lite to microcontrollers with kilobytes of memory, offering portability across Cortex-M devices. Its limitations include manual memory management, restricted operator support, and absence of compiler-based optimizations.

\textbf{CMSIS-NN} \cite{cmsis-nn} delivers ARM-optimized neural network kernels leveraging SIMD instructions. While highly efficient, it requires manual graph construction and lacks high-level training integration or cross-platform support.

\textbf{Glow} \cite{glow} compiles neural networks to optimized machine code using LLVM, achieving strong performance on edge CPUs. However, its compilation times and limited quantization options constrain usability for rapid iteration.

\textbf{MicroTVM} \cite{microtvm} brings Apache TVM's autotuning capabilities to microcontrollers, enabling cross-compilation for diverse targets. Challenges include steep learning curves, fragmented tooling, and incomplete production features (telemetry, privacy).

\subsection{Model Compression Techniques}

\textbf{Quantization} reduces numerical precision from FP32 to INT8 or lower, cutting memory and computation costs. Post-training quantization \cite{ptq} requires no retraining but may degrade accuracy. Quantization-aware training (QAT) \cite{qat} incorporates quantization effects during training, recovering accuracy at the expense of longer training times. Mixed-precision quantization \cite{mixed-precision} assigns different bit-widths per layer, balancing accuracy and efficiency.

\textbf{Pruning} removes redundant connections (unstructured) or entire channels (structured). Structured pruning \cite{structured-pruning} aligns with hardware constraints, avoiding sparse kernels' overhead. Iterative magnitude pruning \cite{lottery-ticket} identifies critical subnetworks while maintaining performance.

\textbf{Knowledge Distillation} \cite{distillation} trains compact student models to mimic larger teachers, transferring learned representations without architectural constraints. Recent variants include attention transfer \cite{attention-transfer} and self-distillation \cite{self-distillation}.

\textbf{Neural Architecture Search (NAS)} \cite{nas} automates model design for target hardware, discovering efficient architectures like MobileNetV3 \cite{mobilenetv3} and EfficientNet \cite{efficientnet}. However, search costs and generalization to custom tasks remain challenging.

\subsection{Intermediate Representations and Compilers}

\textbf{ONNX} \cite{onnx} provides a vendor-neutral model exchange format, facilitating interoperability across frameworks. Its graph representation supports optimization passes but lacks embedded-specific features.

\textbf{MLIR} \cite{mlir} offers a modular compiler infrastructure with customizable dialects, enabling domain-specific optimizations. Projects like IREE \cite{iree} demonstrate MLIR's potential for edge deployment, though adoption requires significant engineering investment.

\textbf{XLA} \cite{xla} optimizes TensorFlow and JAX graphs with fusion, layout optimization, and ahead-of-time compilation. Its Edge TPU backend accelerates inference on Google Coral devices but limits portability.

\subsection{Production Deployment Considerations}

While research frameworks prioritize performance metrics (latency, accuracy), production systems require:
\begin{itemize}
\item \textbf{Privacy}: On-device processing avoids data exfiltration \cite{federated-learning}
\item \textbf{Monitoring}: Drift detection and accuracy canaries ensure reliability \cite{ml-monitoring}
\item \textbf{Compliance}: Auditable policies for regulated industries (healthcare, finance)
\item \textbf{Lifecycle management}: OTA updates, rollback mechanisms, A/B testing
\end{itemize}

Existing frameworks address these concerns piecemeal, creating integration burdens for practitioners.

\subsection{Positioning of Malak Platform}

Malak Platform distinguishes itself through:
\begin{enumerate}
\item \textbf{End-to-end integration}: Unified pipeline from PyTorch training to embedded deployment, eliminating toolchain fragmentation
\item \textbf{Multi-compiler backend}: Supports TVM, MLIR, and XLA simultaneously, enabling hardware-specific optimization without vendor lock-in
\item \textbf{Production-first design}: Built-in telemetry, drift detection, privacy policies, and CLI automation
\item \textbf{Multi-language architecture}: Python for productivity, C++ for portability, Mojo for peak performancebalancing expressiveness and efficiency
\item \textbf{Validated applications}: Three real-world use cases (vision, energy, medical) demonstrating generalizability
\end{enumerate}

Table \ref{tab:comparison} compares Malak Platform against representative frameworks across key dimensions.

\begin{table}[t]
\centering
\caption{Comparison of Edge AI Frameworks}
\label{tab:comparison}
\resizebox{\columnwidth}{!}{%
\begin{tabular}{lccccc}
\hline
\textbf{Feature} & \textbf{TFLite} & \textbf{uTVM} & \textbf{CMSIS-NN} & \textbf{Glow} & \textbf{Malak} \\
\hline
Training Integration & \checkmark & \texttimes & \texttimes & \texttimes & \checkmark \\
Multi-Compiler & \texttimes & \checkmark & \texttimes & \texttimes & \checkmark \\
Hardware Abstraction & Partial & \checkmark & \texttimes & \checkmark & \checkmark \\
Quantization (INT4/8) & INT8 & \checkmark & \checkmark & Partial & \checkmark \\
Production Features & Partial & \texttimes & \texttimes & \texttimes & \checkmark \\
Telemetry/Monitoring & \texttimes & \texttimes & \texttimes & \texttimes & \checkmark \\
Privacy Policies & \texttimes & \texttimes & \texttimes & \texttimes & \checkmark \\
\hline
\end{tabular}%
}
\end{table}
